\section{Introduction}

The standard use of a redistricting ensemble is frequentist: an expert generates
$N$ plans from the ReCom Markov chain, computes a summary statistic (compactness,
partisan share) for each plan, and reports the enacted plan's percentile position.
A plan at the 0.3rd percentile is described as ``more compact than 99.7\% of valid plans.''

This framing has two limitations.
First, the percentile is a point estimate with sampling uncertainty (G.4 establishes
that a 1,000-step NC ensemble has ESS $\approx 70$, giving substantial uncertainty
on any tail percentile).
Second, the court wants to know: \emph{how probable is it that this plan was
intentionally chosen for a specific outcome?}
A frequentist percentile does not directly answer this question.

The Bayesian framework answers it directly.
The prior distribution over redistricting plans is the ensemble.
The observation is the enacted plan's position in the ensemble.
The posterior is the updated probability distribution over how extreme the plan
truly is, accounting for the ensemble's finite size and ESS.

The Bayesian Detection Score (BDS) is a summary of this posterior: the probability
that the plan's true compactness percentile is below 5\% (i.e., in the most
compact 5\% of all valid plans), given the ensemble evidence.
BDS $> 0.95$ means there is strong Bayesian evidence the plan is genuinely extreme;
BDS $< 0.20$ means the plan is consistent with a random neutral draw.

This paper develops the BDS formally and applies it to five study states
(NC, WI, GA, PA, TX) using the G.1 GerryChain ensemble data.
